% Chapter heads: a small gray line "Chapter N", the title below, and a thin gray rule under it.
% KOMA-Script's own interface: \chapterformat is used for numbered chapters only, so \chapter*,
% \addchap, and the bibliography get the title (and the rule) alone.
\setkomafont{chapter}{\normalfont\huge\bfseries}
\renewcommand*{\raggedchapter}{\raggedright}
\renewcommand*{\chapterformat}{{\normalfont\large\color{black!60}\chapapp~\thechapter}}
\makeatletter
% The title is not hyphenated: it is set on as many lines as needed.
\renewcommand*{\chapterlinesformat}[3]{%
  \Ifstr{#1}{chapter}{%
    \Ifstr{#2}{}{}{#2\par\nobreak\vspace{1ex}}%
    {\hyphenpenalty=\@M\exhyphenpenalty=\@M #3}%
    \par\nobreak\vspace{.6ex}{\color{black!30}\rule{\linewidth}{0.4pt}}%
  }{\@hangfrom{#2}{#3}}}
\makeatother
% Alternative: the number in a gray box with the rotated "Chapter" next to it, the title right-aligned below.
% To use it, comment out the lines above and uncomment the following ones.
% \setkomafont{chapter}{\normalfont\Huge\scshape\color{black}}
% \renewcommand*{\raggedchapter}{\raggedleft}
% \renewcommand*{\chapterformat}{%
%   \rotatebox{90}{\makebox[2cm]{\resizebox{2.5cm}{!}{\normalfont\Large\scshape\ \chapapp}}}%
%   \resizebox{!}{2.5cm}{\fboxsep=6pt\colorbox{black!40}{\fontsize{85}{100}\selectfont\sffamily\color{white}\thechapter}}}
% \makeatletter
% \renewcommand*{\chapterlinesformat}[3]{%
%   \Ifstr{#1}{chapter}{%
%     \Ifstr{#2}{}{}{#2\par\nobreak}%
%     {\hyphenpenalty=\@M\exhyphenpenalty=\@M #3}%
%     \par\nobreak\vspace{.6ex}{\color{black!30}\rule{\linewidth}{0.4pt}}%
%   }{\@hangfrom{#2}{#3}}}
% \makeatother
